% ============================================================
%  PHASE 5: SYSTEM ARCHITECTURE & DEPLOYMENT
% ============================================================

\section{Phase 5: System Architecture and Deployment}
\label{sec:architecture}\label{sec:phase5}

The Phase~3--4 model is deployed as a three-tier web application with an
optional language-model layer and an institutional portal for counsellors.
The classifier itself is unchanged: \texttt{predictor.py}, the 88-feature
schema and the fitted scaler statistics are used verbatim.

\subsection{Three-Tier Architecture}
\begin{itemize}
  \item \textbf{Frontend:} Next.js (React, App Router) with Tailwind CSS and
        Recharts. A public landing page, the student assessment wizard and the
        counsellor portal are routes of one application.
  \item \textbf{Backend:} FastAPI (Python). It loads the XGBoost model and label
        encoder once at startup and exposes a REST API. Pydantic schemas
        validate every request before it reaches the model.
  \item \textbf{Database:} PostgreSQL via SQLAlchemy (SQLite for local
        development). Predictions, demographics, raw item responses and
        counsellor notes are stored per session, with the probability tables
        stored as JSON columns to support cohort-level aggregation.
\end{itemize}
The services are containerised with Docker Compose (PostgreSQL, FastAPI,
Next.js).

\subsection{Request Flow and the Gemini Layer}
Figure~\ref{fig:flow} shows the end-to-end flow. The XGBoost model remains the
sole source of the degree-family probabilities. A hosted language model
(a hosted Gemini Flash model, called through the Google GenAI API) is inserted between the
model output and the results page for two purposes that the classifier cannot
serve: generating diagnostic questions and attaching Pakistan-specific market
context.

\begin{figure}[h]
\centering
\begin{tikzpicture}[
  node distance=5mm and 7mm,
  box/.style={draw, rounded corners, align=center, minimum height=9mm,
              minimum width=21mm, font=\footnotesize},
  llm/.style={box, fill=blue!8},
  arr/.style={-{Stealth[length=2mm]}, thick}]
  \node[box] (a) {Student\\answers};
  \node[box, right=of a] (b) {XGBoost\\7 probabilities};
  \node[llm, right=of b] (c) {Stage 1\\follow-up\\questions};
  \node[llm, right=of c] (d) {Stage 2\\validation +\\market data};
  \node[box, right=of d] (e) {Results +\\PDF report};
  \draw[arr] (a) -- (b);
  \draw[arr] (b) -- (c);
  \draw[arr] (c) -- (d);
  \draw[arr] (d) -- (e);
\end{tikzpicture}
\caption{Request flow. Shaded blocks are the language-model layer; the
         classifier output is passed through read-only.}
\label{fig:flow}
\end{figure}

\paragraph{Stage 1: targeted follow-up questions.}
The model receives the ranked probabilities, the Holland code and the houses
whose probabilities lie close to the top match, and returns 3--5
scenario-based multiple-choice questions. Each option is tagged with the
degree family it signals. The motivation comes from Table~\ref{tab:recall}: the classifier
confuses houses with overlapping interest profiles, and fixed RIASEC items
cannot be changed because the trees were fitted to their exact wording.
Follow-up questions probe those ambiguous houses without altering the
model's inputs.

\paragraph{Stage 2: validation and market enrichment.}
The model receives the top-ranked houses and the follow-up answers and
returns (i)~an adjusted confidence for each house, (ii)~a short narrative per
top house and (iii)~market indicators (typical starting salary, common roles,
major employers, approximate openings). Stage~2 is configured to use search
grounding so that market figures come from retrieved web sources rather than
from the model's memory.

\paragraph{Design constraints.}
\begin{itemize}
  \item The classifier is never retrained or modified by the LLM layer. The
        XGBoost probability is always displayed next to any adjusted value;
        the adjustment is a presentation-layer signal only, bounded to
        0--100\% and to at most two rank positions.
  \item Every LLM response is tagged with its provenance
        (\texttt{gemini}, \texttt{gemini\_\allowbreak grounded} or
        \texttt{static\_\allowbreak fallback}), and the interface labels reference data
        as such. This prevents static figures from being presented as live
        results.
  \item The layer degrades gracefully. If the API is rate-limited, slow or
        unavailable, the system falls back to a curated question bank and to
        the static benchmarks of Section~\ref{sec:eda-pk}, so the assessment
        always completes.
  \item Market data is identical for all students in a house, so grounded
        lookups are cached per house rather than repeated per student. This
        keeps cost and rate-limit exposure low.
  \item Only derived values are sent to the external API: probabilities, the
        Holland code and multiple-choice answers. No names, contact details
        or free text are transmitted.
\end{itemize}

\subsection{Counsellor Portal}
Career counselling is rare in Pakistani colleges (Section~1), and a counsellor
who does meet a student has no structured starting point. The portal turns
each completed assessment into a case file. A student receives a short
session identifier and shares it with the counsellor, or links the session to
an institution through an invite code. The portal then provides:
\begin{itemize}
  \item \textbf{Cohort overview:} the number of sessions, mean top-1 fit,
        most common Holland codes and a filterable, searchable session table
        for a school, college or counselling centre.
  \item \textbf{Session detail:} the Holland code and differentiation score,
        RIASEC and Big Five profiles, the full seven-house probability
        breakdown and the market summary, so the counsellor can start the
        conversation from evidence instead of a blank page.
  \item \textbf{Private notes:} counsellor notes saved against the session
        and kept separate from what the student sees.
  \item \textbf{Cohort analytics:} house distribution, Holland-code
        frequencies and urban/rural breakdowns, useful to institutions
        planning guidance programmes.
  \item \textbf{Parents' report:} a one-click PDF of the parents' summary
        for the family meeting.
\end{itemize}
Three demo institutions (a university, a school and a counselling centre) are
seeded. The portal reduces preparation time for the counsellor; it is not a
replacement for professional judgement, and its output is framed as a
conversation starter (Section~\ref{sec:limitations}).

\subsection{API Summary}
\begin{table}[h]
\centering
\small
\caption{REST endpoints of the FastAPI backend.}
\label{tab:api}
\begin{tabular}{lll}
\toprule
\textbf{Method} & \textbf{Route} & \textbf{Purpose} \\
\midrule
POST  & \texttt{/api/predict}                   & Inference and session save \\
POST  & \texttt{/api/gemini/followup-questions} & Stage~1 questions \\
POST  & \texttt{/api/gemini/validate-enrich}    & Stage~2 validation and market data \\
GET   & \texttt{/api/sessions}                  & List sessions (org filter) \\
GET   & \texttt{/api/sessions/\{id\}}           & One session \\
PATCH & \texttt{/api/sessions/\{id\}/notes}     & Counsellor notes \\
GET   & \texttt{/api/sessions/analytics}        & Cohort statistics \\
GET   & \texttt{/api/sessions/\{id\}/pdf}       & Parents' PDF report \\
GET   & \texttt{/api/orgs/\{code\}}             & Resolve invite code \\
\bottomrule
\end{tabular}
\end{table}

\subsection{Assessment Wizard}
The six-step wizard (welcome, background, 48 RIASEC items, TIPI, vocabulary,
results) keeps its state client-side and persists it across page reloads.
After the vocabulary step, the 88-feature vector is assembled server-side
from the raw answers by the original feature-engineering code, which avoids
any training-serving skew. An optional follow-up step sits between
prediction and results.

\subsection{Item Wording and Model Validity}
Because 50.7\% of the total split gain comes from the 48 raw RIASEC items,
their wording cannot be localised without invalidating the trained split
thresholds. The interface therefore adds explanatory hints and a note that
the Investigative items skew towards biology, which can understate
Computer Science fit for technology-oriented students, while leaving the item
text unchanged. Localised items would require re-collecting Pakistani
training data (Section~\ref{sec:nextsteps}).
